\documentclass[11pt]{article}

\usepackage[letterpaper,margin=1in]{geometry}
\usepackage{microtype}
\usepackage[dvipsnames]{xcolor}
\usepackage{enumitem}
\usepackage[most]{tcolorbox}
\usepackage[hidelinks]{hyperref}

\emergencystretch=3em

\definecolor{infoblue}{HTML}{2C3E50}
\definecolor{rule}{HTML}{BDC3C7}

\newtcolorbox{plain}{
  enhanced, breakable,
  colback=Sepia!7, colframe=Sepia!55, boxrule=0.4pt, arc=2pt,
  left=7pt, right=7pt, top=4pt, bottom=4pt,
  fontupper=\small,
  coltitle=Sepia!72!black, fonttitle=\bfseries\sffamily\footnotesize,
  title={THE WHOLE MODEL}
}

\newcommand{\code}[1]{\texttt{#1}}

\title{\bfseries A Practical Internet-Reachability Test\ for FedRAMP}
\author{Matthew Venne \\ \small Chief Technology Officer, stackArmor}
\date{July 2026}

\begin{document}
\maketitle

\begin{center}
\itshape A plain-language companion to ``Internet Reachability at Scale under
the FedRAMP VDR/VER Rules.''
\end{center}

\bigskip

\section*{The short version}
FedRAMP asks providers to decide, for every vulnerability finding, whether a
payload from the public internet might trigger it. That is broader than asking
whether the affected server has a public IP. A crafted request can pass through
an application and reach a private database, parser, queue worker, or logging
library.

The definition is technically sound. The operational problem is that ordinary
scanner findings usually do not say whether a CVE can be triggered after its
input crosses a component boundary. At scale, a provider cannot manually
research that question for every finding on every asset.

The paper recommends a current triage rule:

\begin{plain}
A finding is internet-reachable when its vulnerable network service can be
triggered directly, or when a known indirect trigger reaches the component and
current assessment and continuous-monitoring evidence does not show that the
relevant trigger class is prevented before vulnerable processing.
\end{plain}

That rule is intentionally practical. It uses information providers already
have, promotes exceptional indirect triggers once per CVE, and keeps a more
detailed application-wide data-flow model optional.

\section*{Why the scanner cannot finish the job}
A scanner can usually identify the affected package, CVE, version, and CVSS
vector. It may also report the listening service and a short vulnerability
description. That is enough to automate much of the direct-exposure question.

The missing fact is the trigger path. A local parser flaw may be completely
unrelated to internet input, or it may process every uploaded document. A
logging library may have no listener at all but still receive attacker-controlled
strings from a public API. CVSS attack vector helps describe the vulnerable
component; it does not reliably say whether another component can carry the
trigger to it.

This is primarily an upstream metadata problem. Vendors, CNAs, vulnerability
publishers, and scanner feeds should distribute a reusable field such as
\code{indirect\_internet\_trigger}, with provenance and enough trigger detail to
support review. Today they generally do not.

\section*{Three checks, in order}
The operational method asks three questions.

\begin{enumerate}[leftmargin=1.8em]
\item \textbf{Can the internet directly deliver what this flaw requires?}
  Match public routing, live listeners, protocol, version, and relevant edge
  controls to the vulnerability's direct trigger. A public host does not make
  every installed package reachable.

\item \textbf{Is this a known indirect-trigger vulnerability?}
  Record a reusable CVE-level status: confirmed, ruled out, or unknown. A vendor
  advisory, threat-intelligence report, exploit analysis, reviewed agentic
  proposal, or analyst decision can confirm the status once and reuse it across
  the fleet.

\item \textbf{Does the component receive internet-derived content, and is the
  trigger currently prevented?}
  A coarse component archetype can answer the first half. Current assessment,
  DAST, remediation, and retest evidence can support the second half for the
  trigger or input-control classes they actually cover.
\end{enumerate}

The paper calls that last prevention assertion $P_0$. If $P_0$ is current and
in scope, a known downstream trigger can be classified NIRV. If the evidence is
missing, stale, contradicted, or unrelated to the trigger, $P_0$ is zero and the
finding remains IRV.

\section*{The evidence already exists, but it has limits}
For services following the Continuous Monitoring Playbook, the
\href{https://www.fedramp.gov/resources/documents/Continuous_Monitoring_Playbook.pdf}{FedRAMP
Continuous Monitoring Playbook} requires recurring web-application scanning,
coverage of web interfaces or an approved sample, individual tracking of
findings, and assessor-validated scanner configuration. Providers maintain that
evidence and remediate findings. Assessors review its scope and integrity and
verify provider-run scan results.

That division of labor matters. The provider produces and refreshes the
evidence. The 3PAO collects or reviews it during assessment and validates the
method. The model does not ask the assessor to invent a complete application
data-flow graph or research every CVE on the provider's behalf.

A current, properly scoped DAST program is meaningful sanitation and validation
evidence. It is also not universal proof. Automated testing can miss authenticated
paths, unusual formats, new trigger techniques, and code that the test plan never
exercises. The
\href{https://owasp.org/www-project-web-security-testing-guide/v42/2-Introduction/README}{OWASP
Web Security Testing Guide} makes the same point: automated scanning belongs in a
balanced testing program.

The practical line is straightforward:

\begin{itemize}[leftmargin=1.5em]
\item A clean, current scan with validated scope, relevant coverage, remediated
  findings, and successful retest can support $P_0$ for the tested class.
\item A generic clean scan with unknown scope cannot.
\item A new relevant finding, failed retest, material application change,
  uncovered interface, alternate path, or new trigger technique revokes or
  narrows $P_0$.
\end{itemize}

\section*{Why reachability converges toward accessibility}
The model makes two explicit triage presumptions.

First, a CVE whose indirect-trigger status is unknown is not automatically
promoted across every component that ever handles internet-derived data. Doing
that would make nearly every connected system IRV and erase the prioritization
the field is meant to provide. The unknown status remains visible and is queued
for risk-ranked enrichment.

Second, a current $P_0$ assertion is trusted within its validated scope until
contrary evidence or a material change revokes it. This reuses the assessment and
ConMon work providers already perform.

The resulting operational IRV list therefore converges toward the directly
internet-accessible surface, plus known indirect triggers that are not covered by
current prevention evidence. Reachability and accessibility do not mean the same
thing. They produce similar operational lists because tested prevention closes
many downstream trigger paths, while the exceptional paths are promoted
explicitly.

That is a conditional convergence, not a promise that downstream content is
always safe. Its residual risks are equally explicit: missing upstream trigger
metadata and incomplete or stale prevention evidence.

\section*{Three examples}
\begin{itemize}[leftmargin=1.5em]
\item \textbf{SQL construction behind an application tier.}
  Internet-derived values reach a private database. The validated DAST program
  covers SQL injection on the relevant interface, relevant findings were
  remediated and retested, and no bypass or material change is known. The current
  $P_0$ assertion makes the database finding NIRV. A new SQL-injection finding or
  uncovered interface resets $P_0$, and the finding returns to IRV.

\item \textbf{A logging-expression trigger.}
  Threat intelligence confirms that a crafted string can cross process
  boundaries and trigger the affected logging library. The component receives
  internet-derived log content. Ordinary SQL-injection and cross-site-scripting
  DAST does not cover this mechanism, so $P_0$ is zero and the finding is IRV.

\item \textbf{An administrative service behind IAP or VPN.}
  A separate, bypass-free edge authorizes the connection before traffic reaches
  the backend and is the backend's only ingress path. That can remove the
  backend's direct internet path. A login implemented by the backend application
  does not do the same thing because the request has already reached the
  component; application authentication instead informs exploitability or an
  attested mitigation.
\end{itemize}

\section*{What happens when a new trigger appears}
Log4Shell is the useful pattern. Before the trigger was understood, ordinary
scanner output and web testing did not say that a string written to a log could
activate a private library. Once the mechanism became known, the CVE-level
indirect-trigger status could be promoted across every affected deployment
archetype that received attacker-controlled log content.

That promotion is why trigger metadata belongs upstream. Research it once,
publish it with provenance, and let every scanner and vulnerability-management
platform reuse it. Agentic analysis can propose missing profiles, but it should
cite sources, preserve unknowns, and require stronger review before concluding
that no indirect trigger exists. An agent cannot recover a technical fact that
no authoritative source disclosed.

When new intelligence shows that an existing prevention assertion does not cover
the trigger, $P_0$ is revoked and the affected findings are re-evaluated. The
control is not allowed to remain true merely because last month's scan was clean.

\section*{The optional future state is not today's requirement}
A richer system could import or derive a deployment graph, record how content
moves between components, attach transformations such as parameterized queries
to each path, and bind tests to the deployed artifact. That would improve
precision and make stronger NIRV claims portable between tools.

The industry does not yet have a broadly adopted way to combine those artifacts
into an independently verifiable, vulnerability-specific reachability record
that vulnerability-management platforms can consume. The full paper includes a
non-normative equation showing how such a system could work. It is an
interoperability direction, not a prerequisite for the current method and not a
claim that providers must build a complete application-wide graph now.

\section*{A brief word about WAFs}
FedRAMP allows a verified perimeter control that blocks the triggering payload
before vulnerable processing to remove IRV status. The paper recommends a more
durable operational posture for volatile WAF rules: keep the intrinsic
reachability visible and attach a signed mitigation or VEX assertion describing
the rule, its scope, and its evidence. If the rule is bypassed, disabled, or
routed around, the record fails visibly instead of leaving behind a stale NIRV
classification. KEV remediation dates continue to apply even when a vulnerability
is fully mitigated.

\section*{The bottom line}
\begin{enumerate}[leftmargin=1.8em]
\item Reachability is evaluated per finding, not assigned once to an asset.
\item Direct exposure can be automated from network, listener, protocol, and
  runtime evidence.
\item Exceptional indirect triggers should be enriched once per CVE and reused
  across the fleet; today's scanners rarely provide that metadata.
\item Existing assessment, DAST, remediation, and retest evidence supports a
  scoped, revocable prevention assertion without requiring a complete data-flow
  graph.
\item Under those declared presumptions, the operational IRV list converges
  toward direct accessibility plus known uncovered indirect triggers.
\end{enumerate}

The method is not exhaustive exploit-path proof. It is a scalable triage model
that states exactly what it assumes, preserves the unknowns, and identifies the
upstream metadata improvement that would make the result more precise.

\smallskip
\noindent\textcolor{rule}{\rule{\textwidth}{0.4pt}}
\smallskip

\noindent\small\emph{Read the
\href{https://stackarmor.github.io/rfc-fedramp-vdr/internet-reachability.pdf}{full
technical paper}.}

\end{document}
